\section{Preliminaries}
\label{sec:prelims}

In this section, we introduce definitions and notation that we will use in the remainder of the paper.

\paragraph{Boundary-linked expander decomposition.}
In this paper we work heavily with boundary-linked EDs~\cite{boundlink}. A boundary-linked ED is the same as a classical ED except that Property~\eqref{property:expanderclassic} of \cref{def:expdecclassic} is strengthened. 

For a cut $S$ in a cluster $U$, the boundary $\unbor{S}{U}$ of $S$ with respect to $U$ is the number of edges that go from $S$ to the outside of $U$, i.e. $\unbor{S}{U}=|E(S,V \setminus U)|$. For $U \subseteq V$ and $\tau \geq 0$, the $\tau$-boundary linked subgraph of $G$ on $U$, denoted by $G[U]^{\tau}$, is the subgraph of $G$ induced by $U$ with additional \smash{$\tau \cdot \unbor{\{u\}}{U}$} self-loops attached to every $u \in U$.
	For cuts $S \subseteq U$, we adopt the following shorthand notation:
	\begin{itemize}
		\item the volume $\vol_{G[U]^{\tau}}(S)$ is denoted by \smash{$\vol_U^{\tau}(S)$};
		\item the sparsity $\Phi_{G[U]^{\tau}}(S)$ in $G[U]^{\tau}$ is denoted by $\Phi_U^{\tau}(S)$, which is equal to
		\begin{equation*}
			\Phi_U^{\tau}(S) = \frac{\uncut{S}{U}}{\min\{\vol_U^{\tau}(S), \vol_U^{\tau}(U \setminus S)\}} \, .
		\end{equation*}
	\end{itemize}
Let $\phi \in (0,1), b \in [\phi, 1)$ be some parameters. Most of the time in this paper, the value of~$\tau$ for all instances of boundary linked subgraphs will be the same and equal to the ratio $b/\phi$. Hence we adopt additional shorthand notation: we use \smash{$G[U]^{\circ}$}, \smash{$\vol_U^\circ(S)$}, \smash{$\Phi_U^{\circ}(S)$}, \smash{$\Phi_U^{\circ}$} for referring to \smash{$G[U]^{{b/\phi}}$, $\vol_U^{b/\phi}(S)$, $\Phi_U^{b/\phi}(S)$, $\Phi_U^{b/\phi}$} respectively.
One can observe that
\begin{equation*}
	\vol_U^\circ(S) = \vol_U^{b/\phi}(S) = \vol(S)+\left(\frac{b}{\phi}-1\right)\unbor{S}{U} \, .
\end{equation*}
\noindent
Then a boundary-linked ED can be defined as follows.

\begin{definition}[Boundary-linked expander decomposition~\cite{boundlink}]
	\label{def:expdec}
	Let $G=(V,E)$, let $b, \epsilon, \phi \in \, (0,1)$ be parameters such that $b \ge \phi$, and let $\gamma \ge 1$ be an error parameter. A partition $\mathcal{U}$ of $V$ is a $(b, \epsilon, \phi, \gamma)$-boundary-linked expander decomposition (for short, $(b, \epsilon, \phi, \gamma)$-BLD) of $G$ if
	\begin{enumerate}
		\item \label{property:crossing} $\frac{1}{2}\sum_{U \in \mathcal{U}} \uncutg{U} \le \epsilon |E|$, and
		\item \label{property:expander} for every $U \in \mathcal{U}$, $G[U]^{\circ}$ is a $ \phi/\gamma$-expander.
	\end{enumerate}
\end{definition}

\paragraph{Expander decomposition sequence.} For a partition $\mathcal{U}$ of $V$ (think of $\mathcal{U}$ as an ED of $G$), we denote by $E \setminus \mathcal{U}$ the set of inter-cluster (or crossing) edges with respect to $\mathcal{U}$, i.e.
\begin{equation*}
	E \setminus \mathcal{U} = E \setminus \bigcup_{U \in \mathcal{U}} \binom{U}{2} \, .
\end{equation*}
Analogously, we let $G\setminus \mathcal{U} = (V,E\setminus \mathcal{U})$ be the subgraph of $G$ obtained by removing the intra-cluster edges in $\mathcal{U}$.  For a sequence of partitions $\mathcal{U}_1, \dots, \mathcal{U}_\ell$ of $V$ and $i \in [\ell]$, we define $G^{\remove}_1=G$ and denote by $G^{\remove}_{i+1} = G^{\remove}_{i} \setminus \mathcal{U}_i$ the subgraph of $G$ obtained by removing the intra-cluster edges of the first $i$ partitions.

\begin{definition}[Removal-based ED sequence]
	\label{def:expdecrec}
	Let $G=(V,E)$, let $\epsilon, \phi \in (0,1)$, let $\ell \ge 1$, and let $\mathcal{U}_1, \dots, \mathcal{U}_\ell$ be a sequence of partitions of $V$. The sequence $\mathcal{U}_1, \dots, \mathcal{U}_\ell$ is an $\ell$-level removal-based $(\epsilon, \phi)$-ED sequence (for short, $\ell$-level $(\epsilon, \phi)$-RED or $(\epsilon,\phi,\ell)$-RED) of $G$ if, for all $i \in [\ell]$, $\mathcal{U}_i$ is an $(\epsilon, \phi)$-ED of the graph~$G^{\remove}_i$.
\end{definition}


\paragraph{Sparsifier-specific notation.}
Since our algorithms will be working on sparsifiers of the input graph, it is convenient to have short-hand notation for the above quantities in the sparsifier. For a weighted subgraph $H=(V,E',w)$ of $G$, a cluster $U \subseteq V$ and a cut $S \subseteq U$, we have ``sparsified'' counterparts $\wcutg{S}$, $\wvol{S}$, $\wcut{S}{U},\wbor{S}{U}$ for $\uncutg{S}$, $\vol(S)$, $\uncut{S}{U},\unbor{S}{U}$. More precisely, $\wcutg{S}=w(S,V\setminus S)$, $\wcut{S}{U}=w(S,U\setminus S)$, $\wbor{S}{U}=w(S,V\setminus U)$, and $\wvol{S}$ is the sum of weighted degrees of vertices of $S$ in $H$. Then, $H[U]^\circ$ denotes the subgraph of $H$ induced by $U$ (retaining edge weights from $w$), where every vertex is attached  \smash{$b/\phi \cdot \wbor{\{x\}}{U}$} self-loops. We then also introduce notation for the ``sparsified'' version of $\unvolc{S}{U}$ and $\unspc{S}{U}$: we denote $\vol_{H[U]^{b/\phi}}(S)$ and $\Phi_{H[U]^{b/\phi}}(S)$ by $\wvolc{S}{U}$ and $\wspc{S}{U}$ respectively. For clarity, we note that \smash{$\wvolc{S}{U} = \wvol{S}+\frac{b-\phi}{\phi}\wbor{S}{U}$}, reflecting the fact that \smash{$\vol_U^\circ(S) = \vol(S)+\frac{b-\phi}{\phi}\unbor{S}{U}$}.